\begin{table}[htbp]
\centering
\footnotesize
\caption{Problema G (google/gemini-2.5-flash-lite): valores normalizados (G1) versus kg e m (G2), zero-shot.}
\label{tab:bloco3_G1_vs_G2_zero_shot_flashlite}
\begin{tabular}{llrrrrrrrrrr}
\toprule
 &  & \multicolumn{5}{c}{Acur.\ LLM vs real (\%)} & \multicolumn{3}{c}{Acur.\ métrica (\%)} & \multicolumn{2}{c}{Fidelidade (\%)} \\ \cmidrule(lr){3-7} \cmidrule(lr){8-10} \cmidrule(lr){11-12}
Classes & Atributos & G1 & G2 & $\Delta$ & G1 (inv.\ perm.) & G2 (inv.\ perm.) & G1 & G2 & $\Delta$ & G1 & G2 \\
\midrule
Mulher / Homem & semânticos (peso/altura; peso\_kg/altura\_m) & 35,2 & 84,3 & +49,0 & 64,8 & 84,3 & 10,5 & 91,4 & +81,0 & 71,4 & 92,9 \\
Mulher / Homem & neutros (x1/x2) & 57,1 & 48,6 & $-$8,6 & 57,1 & 51,4 & 74,8 & 65,2 & $-$9,5 & 82,4 & 81,4 \\
A / B & semânticos (peso/altura; peso\_kg/altura\_m) & 18,6 & 46,7 & +28,1 & 81,4 & 54,3 & 16,7 & 22,9 & +6,2 & 93,3 & 70,5 \\
A / B & neutros (x1/x2) & 16,2 & — & — & 83,8 & — & 9,5 & — & — & 90,5 & — \\
\bottomrule
\multicolumn{12}{p{0.95\linewidth}}{\footnotesize Fase A zero-shot, R$^2$, Perceptron, médias das sementes. G1 mostra ao LLM os valores normalizados em $[-1, 1]$; G2 mostra os mesmos pontos em kg e m (regra de três: 45–110 kg, 1,50–1,95 m). A métrica é estimada no espaço normalizado nos dois. Acurácia (vs rótulo real) do LLM e da métrica, contexto; Fidelidade = métrica vs LLM. Acur.\ LLM vs real (inv.\ perm.) = max($a$, 1 $-$ $a$): em zero-shot e com nomes neutros (A/B) a correspondência entre nome e classe é arbitrária, e só o valor invariante à permutação compara as linhas. ``—'' = o LLM colapsou em classe única.} \\
\end{tabular}
\end{table}
